\chapter{Joint versus Sequential Optimization}\label{chap:joint-vs-sequential}

% #############################################################################
% WRITING GUIDE (delete before submission)
%
% THE CHAPTER'S CLAIM -- everything in here serves this one sentence:
%
%   The joint program's in-sample advantage over the sequential overlay is
%   inflated by estimation error in proportion to how little of the sample the
%   risk measure uses. The in-sample gap therefore measures overfitting
%   capacity rather than economic value, and the measures with the strongest
%   in-sample case deliver the least out of sample.
%
% Note the order of that sentence. The MECHANISM is the claim; the inversion
% across risk measures is its consequence. This matters for how you defend the
% chapter: the mechanism is supported by the level sweep (14 levels, one
% measure, slopes with confidence intervals) and by the universe sweep, both of
% which are direct evidence. The cross-measure inversion rests on five points
% and is corroboration. Do not lead with the five-point result and do not let
% the chapter depend on it.
%
% PAGE BUDGET (~16 pages)
%   6.1  In-sample case ................ 2 pages
%   6.2  What survives ................. 2-3 pages
%   6.3  Mechanism ..................... 6 pages   <- the chapter
%   6.4  The inversion ................. 2 pages
%   6.5  What does not repair it ....... 2 pages
%   6.6  When justified ................ 1-2 pages
%
% Note 6.3 and 6.4 are SWAPPED relative to the first draft of this outline:
% mechanism first, inversion as its consequence.
%
% TWO THINGS TO FIX BEFORE WRITING
%
% (a) SIGN CONVENTION. Your draft in the empirical chapter defines
%     g = joint - overlay, so NEGATIVE means joint is better. The code, every
%     figure and every table define gap = overlay - joint, so POSITIVE means
%     joint is better. Adopt the code's convention and fix the prose.
%
% (b) MATERIAL TO MIGRATE from the end of
%     chapters/empirical_analysis_and_baseline_results.tex:
%       - the hypothesis paragraphs  -> Section 6.2
%       - the currency-share arm     -> Section 6.2
%       - the decomposition derivation g^oos = g^is + Delta_epsilon
%                                    -> Section 6.1, as the formal backbone
%     That derivation is good and does real work. It already says what this
%     chapter proves empirically.
% #############################################################################

% TODO: roadmap paragraph, in the style of the other chapter openers. One
% paragraph, naming each section and what it establishes, in order.


\section{The In-Sample Case for Joint Optimization}\label{sec:joint-in-sample-case}

% ---------------------------------------------------------------------------
% THE CLAIM
%   Joint optimization weakly dominates the sequential overlay in sample by
%   construction, and the magnitude of that dominance differs systematically
%   across risk measures -- it is largest for the tail measures.
%
% WHY IT IS TRUE
%   1. The nesting argument, stated as a proposition in Chapter 3 and recalled
%      here. The overlay's solution is feasible for the joint program, and the
%      overlay's second step is a partial minimisation of the same objective:
%
%        f(theta_J) = min_{x,psi} f(x,psi) <= min_psi f(x_O, psi) = f(theta_O)
%
%      Use the derivation already drafted in the empirical chapter.
%   2. Then the decomposition, which is the chapter in one equation:
%
%        g^oos = [f(theta_J) - f(theta_O)] + [eps(theta_J) - eps(theta_O)]
%                \_______ g^is _______/     \____ Delta_eps >= 0 ____/
%
%      where eps(theta) = f^oos(theta) - f(theta) is the generalization gap.
%      The second term is non-negative in expectation because theta_J is
%      selected over a strictly larger set and therefore fits more estimation
%      noise. State this as the chapter's organising identity: the same
%      freedom that wins in sample is what costs out of sample. Everything
%      after this section measures Delta_eps.
%   3. Empirical verification: gap_in < 0 on 8 of 1365 unregularized
%      date-cells (0.6%), all at magnitude ~1e-7. Say explicitly that this is
%      solver tolerance, not violation, and that it doubles as an
%      implementation correctness check.
%   4. The in-sample gap by measure, averaged over the three universes:
%
%        SOES 7.23%   ES 3.71%   SMV 2.13%   MAD 1.91%   MV 1.45%
%
% WHAT THIS SECTION MUST SET UP
%   The expectation that the tail measures are the most promising place to
%   apply joint optimization. State that expectation plainly -- it is the one
%   Section 6.4 overturns, and it has to be on the page to be overturned.
%
% DATA: results/in_sample_vs_out_of_sample/gap_summary.csv (unregularized),
%       gaps.csv for the violation count.
% ---------------------------------------------------------------------------


\section{What Survives Out of Sample}\label{sec:what-survives}

% ---------------------------------------------------------------------------
% THE CLAIM
%   The in-sample advantage does not survive out of sample except in the
%   smallest universe, and it declines monotonically in the number of assets.
%
% WHY IT IS TRUE
%   Out-of-sample relative gap with bootstrap confidence intervals. Significant
%   at top-10 for ES, MAD, MV and SMV; at top-50 only for MV and SMV; at
%   top-100 for nothing; and for SOES nowhere at all.
%
% THE HYPOTHESIS (from the migrated text)
%   The joint program must estimate the asset-currency dependence structure
%   directly; the overlay sidesteps it by sequencing, which is a form of
%   dimensionality reduction. So joint should win where estimation is easy and
%   lose where it is not.
%
% THE OBJECTION THIS SECTION MUST SURVIVE
%   Is the decline driven by the currency share or by the asset count? They
%   move together in the main grid. The K-reduction arm separates them: hold
%   the asset universe fixed and cut the tradable currency block to K = 10, 6
%   (top-50) and K = 10, 6, 3 (top-10). Because s = K/(K + 14N), top-10 at
%   K = 3 lands at s = 0.021 against top-50's 0.020 with 140 securities rather
%   than 700. Keep this to a paragraph and a figure -- it is a robustness check
%   on the hypothesis, not a finding.
%
% FIGURE: results/in_sample_vs_out_of_sample/gap_vs_currency_share.pdf
%   (currently fig:gap-vs-currency-share in the empirical chapter; move the
%   float here rather than duplicating it).
%
% DO NOT interpret the cross-measure differences yet. This section establishes
% that the advantage decays; Section 6.3 explains why; Section 6.4 shows the
% decay is systematic across measures. Keep them apart or the argument blurs.
% ---------------------------------------------------------------------------


\section{Mechanism: The Information Budget}\label{sec:mechanism}

% ---------------------------------------------------------------------------
% THIS IS THE CHAPTER. Six pages. Three converging tests of one mechanism.
%
% THE CLAIM
%   Delta_eps -- the excess generalization gap the joint program pays for its
%   larger feasible set -- grows as the number of observations the risk measure
%   uses falls. The in-sample advantage is inflated by exactly that amount.
%
% Open the section by stating the mechanism in terms of an information budget:
% the joint program estimates N + K parameters; what varies across the three
% tests below is how many observations the objective actually uses to do it.
% Each test varies one thing and holds the others fixed. Say that explicitly
% before the subsections -- it is what makes three results into one argument.
% ---------------------------------------------------------------------------

\subsection{Across Tail Depth Within One Risk Measure}\label{sec:mechanism-level}

% ---------------------------------------------------------------------------
% PUT THIS FIRST. It is the cleanest test and the chapter's strongest evidence:
% it holds the risk measure, the universe, the constraint set and the data
% fixed, and varies only how much of the sample the objective uses.
%
% THE CLAIM
%   Within a single risk measure, looking deeper into the tail raises the
%   in-sample advantage without raising the realized one. The additional
%   in-sample advantage is entirely estimation error.
%
% THE DESIGN
%   ES swept over 14 confidence levels (50, 60, 70, then 2.5-point steps from
%   72.5 to 97.5) on all three universes, 42 cells. Tail observations in the
%   estimation window fall from 1044 at level 60 to 65 at level 97.5 while the
%   number of free weights is unchanged. That ratio IS the information budget,
%   and the sweep varies it by a factor of 16 with nothing else moving.
%
% THE NUMBERS
%   Slope of the relative gap per percentage point of confidence level:
%
%                   in-sample                    out-of-sample (annual)
%     top-10    +0.00162 [0.00135, 0.00186]   +0.00005 [-0.00063, 0.00076]
%     top-50    +0.00100 [0.00078, 0.00124]   +0.00016 [-0.00029, 0.00064]
%     top-100   +0.00060 [0.00038, 0.00087]   +0.00037 [-0.00006, 0.00079]
%
%   In-sample: significantly positive at every universe. Out-of-sample: zero
%   at every universe. At top-10 the in-sample gap runs from 3.11% at level 60
%   to 10.59% at 97.5 while the realized gap sits at 2.2-2.8% throughout.
%
%   The overfit ratio (realized decay relative to 1/N's, which optimizes
%   nothing and so isolates estimation error) rises in the same shape:
%   top-10 from 1.34 at level 60 to 2.01 at 97.5; top-100 from 3.07 to 5.88.
%
% THE SENTENCE TO WRITE
%   The in-sample gap decomposes into a level-independent structural component
%   (~2.5% at top-10, which survives) and a level-increasing estimation
%   component (which does not). Deeper tail targeting buys no additional real
%   advantage; it buys the joint program more room to overfit a thinning tail.
%
% THE OBJECTION THIS MUST SURVIVE -- and it is a real one, so do not bury it.
%   f_out is scored on the quarter actually held, whose tail thins to ~3
%   observations at level 95 and ~2 at 97.5. At that point f_out is an extreme
%   order statistic, not an expectation. A second horizon (the following year,
%   ~252 days) is therefore scored alongside it, lifting the level-97.5 tail
%   to ~6 observations. The quarterly horizon suggests the realized gap
%   DECLINES at high levels; the annual horizon shows it is FLAT and agrees
%   with the pooled production backtests where the quarterly one does not.
%   Report both, say which you rely on and why. Handling this openly is what
%   makes the section credible; an examiner will go looking for it.
%
% DATA: results/in_sample_vs_out_of_sample_es_at_diff_levels/
% ---------------------------------------------------------------------------

\subsection{Across Problem Dimension}\label{sec:mechanism-dimension}

% ---------------------------------------------------------------------------
% THE CLAIM
%   Holding the risk measure fixed and raising the number of assets has the
%   same effect: more parameters against the same observations, so more of the
%   in-sample advantage is estimation error.
%
% THE NUMBERS
%   Overfit ratios for ES rise from 1.38 at top-10 to 3.09 at top-100. Both
%   programs overfit more as the universe grows, and the joint program
%   overfits more than the overlay in 12 of 15 cells (paired, CIs clear of
%   zero). The three exceptions -- top-10 at levels 60, 70, 80 -- are exactly
%   the cells where the joint advantage survives out of sample.
%
% THE SENTENCE TO WRITE
%   The joint program's advantage survives if and only if it does not overfit
%   more than the overlay does. That is a mechanism, not a correlation, and it
%   is worth stating as the section's punchline.
%
% Note that both programs still beat 1/N in absolute realized risk everywhere
% (0.38 to 0.69 of it). Say so -- it prevents the reader concluding that
% optimization does not work, when the claim is narrower.
% ---------------------------------------------------------------------------

\subsection{Across Risk Measures}\label{sec:mechanism-measures}

% ---------------------------------------------------------------------------
% THE CLAIM
%   The same mechanism operates across risk measures. Deviation measures
%   (variance, MAD, semi-variance) use every observation; tail measures (ES,
%   SOES) use only the tail. The tail measures therefore have the largest
%   in-sample gap and the least realizable one.
%
% THE COMPETING EXPLANATION THIS MUST RULE OUT
%   ES and SOES are not location-free -- they carry an implicit mean. Is their
%   behaviour driven by that rather than by tail focus? Test by de-meaning,
%   which makes them location-free. Result: the in-sample ceiling falls but
%   does not vanish. So the implicit mean explains part and the tail focus the
%   rest. Report the magnitude of the fall; this is the experiment that
%   converts a plausible story into a tested one.
%
% DATA: results/in_sample_vs_out_of_sample_demeaned/
% ---------------------------------------------------------------------------


\section{The Inversion Across Risk Measures}\label{sec:the-inversion}

% ---------------------------------------------------------------------------
% THE CLAIM
%   Because the in-sample gap is inflated in proportion to the information
%   deficit, its magnitude predicts the realized benefit NEGATIVELY across
%   risk measures.
%
% This section is now a CONSEQUENCE of Section 6.3, not an independent
% finding. Write it that way: "the mechanism implies an ordering, and the
% ordering is observed." That framing is both more honest and more persuasive
% than presenting it as a standalone discovery.
%
% THE NUMBERS -- averaged over the three universes, unregularized:
%
%     measure    in-sample   out-of-sample   share realized
%     SOES         7.23%        0.47%             7%
%     ES           3.71%        1.14%            31%
%     SMV          2.13%        1.87%            88%
%     MAD          1.91%        1.18%            62%
%     MV           1.45%        1.52%           105%
%
%   Spearman rank correlation between the in-sample and out-of-sample gap
%   across the five measures: -0.50 (top-10), -0.20 (top-50), -0.60
%   (top-100). Negative at every universe size.
%
% THE HONESTY REQUIREMENT -- do not skip this, it is what makes the claim
% survive a viva:
%   Five measures is a thin basis for a rank correlation and the correlation
%   is not itself a powered test. Say so in one sentence, and point to
%   Section 6.3.1: the level sweep reproduces the same effect WITHIN one
%   measure with 14 points and confidence intervals, which is what makes the
%   cross-measure ordering credible rather than a coincidence over five
%   points. Presented in that order the thin result is corroboration of a
%   well-supported mechanism instead of a load-bearing claim.
%
% FIGURE: not yet built. Whatever form you choose, it has to put the
% in-sample magnitude and the realized share in the same coordinate system.
% The existing five-panel figure cannot express this claim because neither
% variable of the claim is on an axis.
% ---------------------------------------------------------------------------


\section{What Does Not Repair It}\label{sec:what-does-not-repair}

% ---------------------------------------------------------------------------
% Defence. Pre-empts the two obvious objections. Two pages, no more.
% ---------------------------------------------------------------------------

\subsection{Regularization}\label{sec:no-repair-regularization}

% ---------------------------------------------------------------------------
% METHODOLOGICAL POINT FIRST -- it justifies the scope of the whole chapter.
%   The in-sample gap framework does NOT extend to penalised programs. The
%   nesting proposition requires the scorer to be the objective. Regularized
%   joint and overlay programs minimise Risk + penalty with penalties selected
%   per date that differ between the two, so they optimise different
%   objectives while being scored on a third. Nothing is bounded and the
%   in-sample/out-of-sample difference is not the decay of anything.
%
%   EVIDENCE THAT THIS IS NOT HYPOTHETICAL: the ceiling gap_in >= 0 is
%   violated on 30.0% of regularized date-cells (MV 44.7%, SMV 38.5%,
%   MAD 20.1%, ES 16.8%) against 0.6% at solver tolerance for unregularized.
%   Cite these numbers. They turn the restriction of this chapter to the
%   unregularized family from an apparent omission into a consequence of the
%   proposition -- which is a much better thing for an examiner to find.
%
% THE CLAIM, answered out of sample only, by comparing realized risk of the
% production backtests (no in-sample quantity, no ceiling required):
%   Regularization reduces the absolute estimation error both programs suffer
%   without changing which of them is lower.
%
%   ES at level 80, relative reduction in realized ES_80, joint vs overlay:
%
%     universe      unregularized     regularized
%     top-10          +3.00% (sig)     +2.01% (sig)
%     top-50          +1.76% (sig)     +1.40% (sig)
%     top-100         +0.15% (ns)      -0.21% (ns)
%
%   Same pattern, same cut-off, slightly smaller under regularization. And the
%   overfit ratio falls sharply: at top-100 from 3.09 to 1.93, with realized
%   risk relative to 1/N improving from 0.384 to 0.350. Regularization pulls
%   both programs toward the same floor without changing their ordering.
% ---------------------------------------------------------------------------

\subsection{A Mean Term in the Objective}\label{sec:no-repair-mean}

% ---------------------------------------------------------------------------
% THE CLAIM
%   Adding a mean term helps the deviation measures substantially and does
%   little for ES and SOES, which already carry the mean implicitly.
%
% Frame as consistent with the mechanism: the measures with information to
% spare can afford to estimate a mean; the information-starved ones cannot.
%
% The nesting proposition fails here too, and for the same reason -- gamma is
% selected per date and the joint and overlay programs routinely select
% different values, so the two minimise different objectives. Say so and rely
% on realized quantities.
%
% DATA: results/in_sample_vs_out_of_sample_w_mean/, top_runs_w_mean.csv
% ---------------------------------------------------------------------------


\section{When Joint Optimization Is Justified}\label{sec:when-justified}

% ---------------------------------------------------------------------------
% THE CLAIM
%   Joint optimization is justified in small universes with variance-based
%   measures, and specifically not with the tail measures that look most
%   promising in sample.
%
% State that this is counterintuitive and that it follows directly from the
% mechanism: the tail measures' in-sample appeal IS their estimation problem.
%
% MAGNITUDE, net of transaction costs (this matters -- the gap analysis
% ignores costs and the production comparison does not, which closes the
% objection that higher currency turnover eats the benefit):
%   At top-10, joint reduces realized volatility by 2.5-5.7% and realized ES
%   by 1.8-3.9% relative to the overlay, with no detectable effect on mean
%   return. At top-100 nothing is significant on any measure.
%
% THE SHARPE QUESTION, and why it is the wrong instrument:
%   Sharpe divides an unestimable numerator by the improved denominator. The
%   implied Sharpe change from a 2.5% volatility reduction is ~0.003 against a
%   detection threshold of ~0.15 at this sample length. Report it as a bound
%   ("we can exclude a Sharpe improvement above roughly 0.10") rather than as
%   a null. Note that the same bootstrap on the same return series DOES detect
%   the volatility and ES effects, so the null is not a dead instrument -- it
%   is a statement about what Sharpe can resolve.
%
% THE HONEST LIMIT -- state it plainly, hedging would weaken the chapter:
%   At the universe sizes relevant for international equity portfolios
%   (700-1400 securities against 14 currencies), the two approaches are
%   statistically indistinguishable out of sample, and the sequential overlay
%   is cheaper to compute. That is the practical conclusion.
% ---------------------------------------------------------------------------
